\documentclass[11pt,a4paper]{article}

\usepackage[margin=2.5cm]{geometry}
\usepackage{amsmath}
\usepackage{hyperref}
\usepackage{listings}
\usepackage{xcolor}

\lstset{
  basicstyle=\ttfamily\small,
  breaklines=true,
  frame=single,
  columns=fullflexible,
  backgroundcolor=\color{gray!5}
}

\title{FAC + ODAT-SE 第一阶段配置搜索流程整理}
\author{}
\date{2026-05-23}

\begin{document}

\maketitle

\section{项目背景}

本项目面向 EBIT 实验光谱中高电荷态重元素离子的 UTA
（unresolved transition array）或 blended peak 归属问题。目标是利用
FAC 计算给定 configuration 与 radial potential 策略下的能级、跃迁能量、
A-value 与角动量信息，再由 ODAT-SE 或其前置的搜索流程反演可能的
configuration 与势优化策略。

FAC 在这里承担顺问题求解器的角色：

\begin{lstlisting}
configuration + radial potential -> FAC -> levels + transition energies + A-values
\end{lstlisting}

ODAT-SE 或简化搜索器承担逆问题搜索层的角色：

\begin{lstlisting}
experimental peak features -> search configuration + potential strategy
\end{lstlisting}

当前阶段不是完整 CRM 谱拟合，而是第一阶段筛选。其目标是判断一组
configuration / potential 是否能在实验峰附近产生合理的理论线团。

一个重要限制是，实验峰通常是 UTA 或 blended peak，因此不能假设

\begin{lstlisting}
one experimental peak = one theoretical transition
\end{lstlisting}

而应采用

\begin{lstlisting}
one experimental peak = one theoretical line group
\end{lstlisting}

此外，FAC 阶段只有 A-value，而真实谱线强度应满足

\[
I_{ul} \propto N_u A_{ul}.
\]

其中 $N_u$ 是上能级 population，需要 CRM 才能给出。因此当前阶段
A-value 只作为 radiative-potential weight，不能当作真实实验强度。

\section{参考 FAC 输入文件}

今天首先阅读了项目中的 \texttt{Bi22.py}，并将其作为当前 FAC 输入脚本的
标准风格。该文件的主要结构如下：

\begin{itemize}
  \item 设置元素和电子数，例如 \texttt{Z = 83}, \texttt{K = 61}；
  \item 使用 \texttt{fac.SetAtom(a)} 设置元素；
  \item 使用 \texttt{fac.Closed(...)} 指定闭壳层；
  \item 使用 \texttt{fac.Config(..., group='nX')} 定义 configuration group；
  \item 使用
    \texttt{ConfigEnergy(0) -> OptimizeRadial(...) -> ConfigEnergy(1)}
    生成并修正单中心势；
  \item 使用 \texttt{Structure} 输出能级；
  \item 使用 \texttt{TransitionTable} 输出跃迁和 A-value；
  \item 可选地计算 CE 和 CI 表。
\end{itemize}

随后阅读了 FAC 手册 \texttt{/home/song/projects/fac-master/doc/manual.tex}。
其中确认了以下关键点：

\begin{itemize}
  \item PFAC，即 Python 输入接口，是推荐使用方式；
  \item \texttt{Config} 支持类似 \texttt{5[s,p,d,f,g]} 的紧凑写法；
  \item \texttt{OptimizeRadial} 应优先使用最低的一个或两个 configuration group；
  \item \texttt{MemENTable} 必须在 verbose \texttt{PrintTable(..., 1)} 之前调用；
  \item \texttt{TransitionTable(fn, low, up, m)} 的第 4 个参数不是
    ``最大 multipole 阶数''。
\end{itemize}

FAC 手册和源码确认的 multipole 约定为：

\begin{lstlisting}
m < 0: electric multipole, e.g. -1 = E1
m > 0: magnetic multipole, e.g. +1 = M1
m = 0 or omitted: sum allowed multipoles
\end{lstlisting}

因此 \texttt{+3} 表示 M3，不表示 E1/E2/E3/M1/M2/M3 全部计算。若希望
汇总到三阶 electric 与 magnetic multipoles，目前采用：

\begin{lstlisting}[language=Python]
fac.SetTransitionMaxE(3)
fac.SetTransitionMaxM(3)
fac.TransitionTable(..., 0)
\end{lstlisting}

\section{目标输入文件设计}

今天建立并多次修改了主输入文件：

\begin{lstlisting}
inputs/target_case.yaml
\end{lstlisting}

该文件不是 FAC 直接输入文件，而是用户编辑的上游配置文件。它包含：

\begin{itemize}
  \item 实验参考峰；
  \item 元素、原子序数、价态和电子数；
  \item FAC 固定设置；
  \item configuration 生成规则；
  \item \texttt{OptimizeRadial} 策略；
  \item 搜索和评分参数。
\end{itemize}

当前示例为 W46+：

\begin{lstlisting}
ion:
  element: "W"
  Z: 74
  charge_state: 46
  K: 28
\end{lstlisting}

即 Ni-like tungsten，28 个束缚电子。

参考峰为：

\begin{lstlisting}
5.69 nm
5.87 nm
\end{lstlisting}

FAC 的 \texttt{.tr} 文件中输出的是跃迁能量，因此需要在波长和能量之间转换：

\[
E[\mathrm{eV}] =
\frac{12398.419843320026}{\lambda[\mathrm{\AA}]}.
\]

反向转换为：

\[
\lambda[\mathrm{\AA}] =
\frac{12398.419843320026}{E[\mathrm{eV}]}.
\]

对应当前两个参考峰：

\begin{align}
5.69~\mathrm{nm} &\approx 217.8984~\mathrm{eV},\\
5.87~\mathrm{nm} &\approx 211.2167~\mathrm{eV}.
\end{align}

\section{Configuration 模板}

用户希望不要在输入文件中逐条写展开后的 FAC configuration，而是写成更接近
物理记号的形式。例如：

\begin{lstlisting}
3p6 3d9 nl, n = 4,5,6; l = 0-4
\end{lstlisting}

因此今天设计了 \texttt{configuration\_space.templates}。当前 W46+ 示例为：

\begin{lstlisting}
configuration_space:
  templates:
    bound:
      - id: "ground_3p6_3d10"
        prefix: "3p6 3d10"
        active: null
        required: true
      - id: "hole_3d9_nl"
        prefix: "3p6 3d9"
        active: "nl"
        n: [4, 5, 6]
        l: [0, 1, 2, 3, 4]
        occupancy: 1
        required: false
      - id: "hole_3p5_3d10_nl"
        prefix: "3p5 3d10"
        active: "nl"
        n: [4, 5, 6]
        l: [0, 1, 2, 3, 4]
        occupancy: 1
        required: false
\end{lstlisting}

生成器会将其展开为 FAC 可读形式。例如：

\begin{lstlisting}[language=Python]
fac.Config('3p6 3d10', group='n2')
fac.Config('3p6 3d9 4[s,p,d,f]', group='n3')
fac.Config('3p6 3d9 5[s,p,d,f,g]', group='n4')
fac.Config('3p6 3d9 6[s,p,d,f,g]', group='n5')
fac.Config('3p5 3d10 4[s,p,d,f]', group='n6')
fac.Config('3p5 3d10 5[s,p,d,f,g]', group='n7')
fac.Config('3p5 3d10 6[s,p,d,f,g]', group='n8')
\end{lstlisting}

今天特别修正了一个物理约束：对于电子轨道，必须满足

\[
l \le n - 1.
\]

因此不存在 \texttt{4g}。当输入为 \texttt{n=4, l=0-4} 时，实际只生成
\texttt{4[s,p,d,f]}。

\section{FAC 输入生成脚本}

今天创建了 FAC 输入生成脚本：

\begin{lstlisting}
scripts/generate_fac_input.py
\end{lstlisting}

该脚本读取 \texttt{inputs/target\_case.yaml}，展开 configuration 模板，
生成类似 \texttt{Bi22.py} 的 PFAC 输入文件。例如：

\begin{lstlisting}
python3 scripts/generate_fac_input.py inputs/target_case.yaml -o generated/W46_phase1.py
\end{lstlisting}

生成的示例文件为：

\begin{lstlisting}
generated/W46_phase1.py
\end{lstlisting}

其核心流程为：

\begin{lstlisting}[language=Python]
fac.SetAtom(a)
SetUTA(0)
fac.Closed(...)
fac.Config(...)

fac.ConfigEnergy(0)
fac.OptimizeRadial([...])
fac.ConfigEnergy(1)
fac.GetPotential(...)

fac.Structure(...)
fac.MemENTable(...)
fac.PrintTable(..., 1)

fac.SetTransitionMaxE(3)
fac.SetTransitionMaxM(3)
fac.TransitionTable(..., 0)
fac.PrintTable(..., 1)
\end{lstlisting}

\section{OptimizeRadial 当前处理}

\texttt{OptimizeRadial} 是后续很复杂的问题，目前暂时固定使用基态：

\begin{lstlisting}
radial_potential:
  optimize_radial_groups: ["ground_3p6_3d10"]
\end{lstlisting}

生成器会将语义 ID 映射成实际 FAC group。例如：

\begin{lstlisting}
ground_3p6_3d10 -> n2
\end{lstlisting}

最终生成：

\begin{lstlisting}[language=Python]
fac.OptimizeRadial(['n2'])
\end{lstlisting}

这只是第一版保守策略。以后需要将 \texttt{OptimizeRadial} 纳入搜索变量，例如：

\begin{itemize}
  \item ground only；
  \item ground + low excited group；
  \item 3d-hole family；
  \item 3p-hole family；
  \item weighted mean configuration；
  \item 不同 potential strategy。
\end{itemize}

\section{Configuration 搜索流程}

用户指出，仅仅将输入 configuration 映射成 FAC 输入文件还不够，需要循环寻找
最佳 configuration。因此今天创建了搜索脚本：

\begin{lstlisting}
scripts/run_configuration_search.py
\end{lstlisting}

当前搜索器还不是完整 ODAT-SE，而是一个可工作的 template-subset grid search。
运行方式：

\begin{lstlisting}
python3 scripts/run_configuration_search.py inputs/target_case.yaml --clean
\end{lstlisting}

搜索器执行以下步骤：

\begin{enumerate}
  \item 保留 \texttt{required: true} 的模板；
  \item 枚举 \texttt{required: false} 模板的开关组合；
  \item 为每个 trial 临时生成 PFAC 脚本；
  \item 在独立目录运行 FAC；
  \item 解析 ASCII \texttt{.tr} 文件；
  \item 对参考峰附近的理论线团评分；
  \item 将结果写入 CSV。
\end{enumerate}

当前 W46+ 示例中，基态模板为 required，两个 \texttt{nl} family 为 optional，
因此第一版搜索尝试四种组合：

\begin{lstlisting}
ground only
ground + 3p5 3d10 nl
ground + 3p6 3d9 nl
ground + both nl families
\end{lstlisting}

搜索结果输出为：

\begin{lstlisting}
runs/search_W46/results.csv
\end{lstlisting}

为了避免大量 configuration 时生成大量 \texttt{.py} 文件，今天修改为默认不保留
每个 trial 的 Python 输入脚本，而是使用临时文件运行后删除。若调试需要，可在
YAML 中设置：

\begin{lstlisting}
search:
  keep_generated_scripts: true
\end{lstlisting}

\section{评分方式：能量与波长两种窗口}

FAC 的 \texttt{.tr} 文件中记录 transition energy，单位为 eV。但实验峰通常以
波长给出。因此今天将评分窗口设计为可选：

\begin{lstlisting}
search:
  scoring:
    window_space: "wavelength"
    energy_window_eV: 8.0
    wavelength_window:
      value: 0.1
      unit: "nm"
\end{lstlisting}

若 \texttt{window\_space = "wavelength"}，则对每个参考峰筛选满足：

\[
|\lambda_i - \lambda_{\mathrm{peak}}| \le 0.1~\mathrm{nm}
\]

的 FAC 跃迁线。每条线的波长由其能量反算得到：

\[
\lambda_i[\mathrm{\AA}] =
\frac{12398.419843320026}{E_i[\mathrm{eV}]}.
\]

若 \texttt{window\_space = "energy"}，则使用：

\[
|E_i - E_{\mathrm{peak}}| \le \Delta E.
\]

例如：

\begin{lstlisting}
window_space: "energy"
energy_window_eV: 8.0
\end{lstlisting}

当前两种方式都保留，用户可在 YAML 中切换。

\section{当前 Loss}

对每个实验参考峰，窗口内所有理论跃迁构成一个 theoretical line group。

当前使用 A-value 或 rate 加权计算线团中心：

\[
E_{\mathrm{grp}} =
\frac{\sum_i E_i A_i}{\sum_i A_i},
\]

\[
\lambda_{\mathrm{grp}} =
\frac{\sum_i \lambda_i A_i}{\sum_i A_i}.
\]

当前简化 loss 为：

\[
L \mathrel{+}= r^2 - \log\left(\sum_i A_i + \epsilon\right),
\]

其中 $r$ 根据评分空间不同取：

\[
r = E_{\mathrm{grp}} - E_{\mathrm{peak}}
\]

或

\[
r = \lambda_{\mathrm{grp}} - \lambda_{\mathrm{peak}}.
\]

若某个参考峰窗口内没有理论线，则加入大的 missing penalty：

\begin{lstlisting}
missing_peak_penalty: 1000000.0
\end{lstlisting}

数值下限为：

\begin{lstlisting}
rate_floor: 1.0e-99
\end{lstlisting}

搜索结果中同时记录能量和波长空间的中心与残差：

\begin{itemize}
  \item \texttt{center\_eV}
  \item \texttt{center\_nm}
  \item \texttt{residual\_eV}
  \item \texttt{residual\_nm}
\end{itemize}

\section{Python 环境}

今天检查了 Python 环境。系统 Python 为：

\begin{lstlisting}
/usr/bin/python3
\end{lstlisting}

该环境可以 import \texttt{pfac} 和 \texttt{yaml}，但最初没有
\texttt{numpy/scipy/pandas}。随后为系统 Python 安装并验证：

\begin{lstlisting}
numpy 2.4.6
scipy 1.17.1
pandas 3.0.3
\end{lstlisting}

同时确认：

\begin{lstlisting}
pfac /home/song/projects/fac-master/_install/local/lib/python3.12/dist-packages/pfac/__init__.py
\end{lstlisting}

\section{新增和修改的文件}

今天主要新增或修改了以下文件：

\begin{itemize}
  \item \texttt{inputs/target\_case.yaml}: 主输入配置文件；
  \item \texttt{scripts/generate\_fac\_input.py}: YAML 到 PFAC 输入脚本生成器；
  \item \texttt{scripts/run\_configuration\_search.py}: configuration 搜索和评分脚本；
  \item \texttt{generated/W46\_phase1.py}: 由当前 YAML 生成的单次 FAC 输入示例；
  \item \texttt{docs/run\_generation.md}: 运行说明；
  \item \texttt{docs/claude\_method\_review.md}: 给 Claude 评价方法用的完整说明；
  \item \texttt{docs/daily\_summary\_fac\_odat\_2026\_05\_23.tex}: 本总结文档。
\end{itemize}

\section{当前已知限制}

当前方法仍是第一阶段骨架，主要限制包括：

\begin{enumerate}
  \item 还不是完整 ODAT-SE，只是 template-subset grid search；
  \item 搜索空间目前只支持 optional template 的开关组合；
  \item \texttt{OptimizeRadial} 暂时固定为基态；
  \item loss 还没有引入 global energy shift 或 wavelength shift；
  \item loss 还没有显式使用实验峰宽、峰面积或不确定度；
  \item 当前窗口是 hard window，未来应考虑 Gaussian soft window；
  \item A-value 加权线团中心可能因缺少 population 而有偏；
  \item 尚未包含 CRM population、cascade 或 ion abundance；
  \item 大 configuration set 下 FAC 计算成本可能很高，尚未实现缓存和剪枝；
  \item 不同 configuration 数量的模型之间需要额外复杂度惩罚，否则大模型天然产生更多线。
\end{enumerate}

\section{后续建议}

下一步建议包括：

\begin{enumerate}
  \item 先运行当前 W46+ 四个 trial，检查 FAC 是否稳定完成；
  \item 解析 \texttt{results.csv}，人工查看两个参考峰附近的候选线团；
  \item 将 hard window 改成基于实验峰宽的 Gaussian soft weighting；
  \item 加入 global shift consistency；
  \item 加入模型复杂度惩罚；
  \item 扩展搜索空间，使 $n$、$l$、configuration family 和
    \texttt{OptimizeRadial} strategy 都能作为搜索变量；
  \item 在少量 top candidates 上做简单 Gaussian convolution 验证；
  \item 最后再进入 CRM 完整验证。
\end{enumerate}

\end{document}
